\documentclass{report}

\usepackage[a4paper,margin=1in,footskip=0.25in]{geometry}
\usepackage{pdflscape}
\usepackage{csvsimple}
\usepackage{booktabs}
\usepackage[none]{hyphenat}
\usepackage[colorlinks=true,allcolors=black,urlcolor=blue]{hyperref}
\usepackage{minted}

\title{Animating Digital Arithmetic}
\author{Will Dargan}

\setcounter{tocdepth}{3}

\begin{document}

\maketitle

\pagenumbering{roman}
\setcounter{page}{2}

\ \newpage

\chapter*{Abstract}
\addcontentsline{toc}{chapter}{Abstract}

\newpage \ \newpage

\tableofcontents
\addcontentsline{toc}{chapter}{Contents}
\newpage

\pagenumbering{arabic}

\chapter*{Introduction}
\addcontentsline{toc}{chapter}{Introduction}
This project revolves around an application that enables the user to dynamically
generate educational animations of simple arithmetic and bitwise logic on binary
numbers.

\section*{Aims}
\addcontentsline{toc}{section}{Aims}

\section*{Objectives}
\addcontentsline{toc}{section}{Objectives}

\chapter{Literature Review}

\section{Underpinning Theory}

\section{Technology Stack}
As will be outlined and explained in much more depth in Section
\ref{software arch}, the final solution will be broken down into three major
components: a GUI; a component for generating scripts that animations may be
generated from; and finally a component for generating the animations
themselves. As the second component is largely language and framework agnostic,
the following paragraphs do not contain a dedicated section for it.
However, there is a section on the choice of programming language, so this
component will be considered there.

\subsection{Programming Language}
There are a small few languages that I have some experience with and that would
be suitable for a large scale project. Each of them is listed below with a
brief description of their characteristics:

\paragraph{C++}
An incredibly fast and versatile compiled language that provides the
programmer with direct access to memory without a garbage collector.
C++ supports both functional and object-oriented programming paradigms.

\paragraph{HTML/CSS/JS}
This suite of languages comprises the tech stack used by all modern browsers to
render webpages and enable interactivity. There are a large number of
third-party frameworks for these languages to simplify and speed up development,
including the ability to write HTML/CSS/JS for stand-alone desktop applications.
They can also be used to create a front-end for a web-server written in any of
the other general-purpose languages listed here.

\paragraph{Java}
A very widely used, object-oriented first, language intended to be able to be
run on any hardware that has a Java Virtual Machine (JVM). Java is both a
compiled \textit{and} interpreted language, as source code is first compiled
into an intermediary language that is then interpreted by the JVM at runtime.
Despite this, Java is still highly performant.

\paragraph{Python}
A beginner friendly language supporting both functional and object-oriented
programming paradigms. Python is an interpreted language, making it
considerably slower than some of the other languages listed above. However, it
still sees considerable use in the AI and Machine Learning industry. Python is
also commonly used as a first language for teaching programming in secondary
schools \cite{hadwenbennett2024programming} and universities.
\vspace{3.25ex}

I decided against HTML/CSS/JS as I want my final solution to be a native
application, for which one of the other languages would be a better fit. I
also ruled out C++, as I would rather not put myself through the extra
difficulty programming in it when the extra features offered are
unnecessary for this project.

Of the two remaining languages, both of them require an interpreter to be
installed on the end user's system\footnote{\label{pyinstaller}Python
\textit{can} be bundled with an interpreter into a single executable.}.
As mentioned earlier, Python is the most commonly used language for teaching in
schools \cite{hadwenbennett2024programming}, so with this project positioned as
a teaching aid, I can expect the vast majority of end users to already have a
Python installation on their systems. Most distributions of macOS and Linux
also come with an installation of Python included. Finally, between the two
remaining languages, I have far more experience writing Python than I do Java.
For these reasons I have chosen to use Python for this project.

\subsection{GUI}
There are many GUI libraries on offer for Python, each one with it's own
specialities and capabilities. Below are three of the most common libraries
for writing native desktop GUI applications, all of which are widget toolkits:

\paragraph{PySide6/PyQt6 (Qt)}
These are two separate but very similar libraries containing Python bindings
for the Qt framework, which supports Windows, macOS and Linux, as well as
Android for mobile. Qt is a flexible framework that can be used to create
native-looking GUIs as well as 3D UI elements \& graphics and graphs. Qt also
offers direct integration with databases via SQL.

\paragraph{Tkinter (Tcl/Tk)}
A module in the Python standard library that provides bindings for the Tcl/Tk
GUI framework. Tk supports Windows, macOS and Linux with native-looking GUI
applications. Widgets may be styled with `themed Tk' introduced in Tk 8.5, the
Python bindings for which are supplied in the \mintinline{python}{tkinter.ttk}
submodule.

\paragraph{wxPython}
A library for creating truly native cross-platform GUI applications with
support for Windows, macOS and Linux. wxPython wraps platform specific widgets
in platform agnostic bindings so that the same code may generate truly native
applications for any operating system. This gives rise to some quirks and differences
between platforms, where some options are unsupported. For example, macOS
buttons may not be resized vertically.
\vspace{3.25ex}

While researching Python GUI libraries I came across two other libraries which
were unsuitable for this project, but are notable nonetheless: Kivy, which is a
mobile first framework for writing multi-touch capable applications; and
PyGObject, which is a native GUI library specifically for GNOME based Linux
desktop environments.
\vspace{3.25ex}

I decided not to go with wxPython as the OS specific differences would make it
difficult to develop an application that has a consistent appearance across all
platforms. Both of the remaining frameworks are capable enough for my needs in
this project. My decision was therefore based on the library that I have more
experience with, which in this case is Tkinter.

\subsection{Animations}
\paragraph{Manim Community}
\paragraph{OpenGL}
\paragraph{pygame-ce}

\chapter{Design}

\section{Software Architecture} \label{software arch}

\begin{landscape}

	\section{Requirements}

	\subsection{Essential Features}
	\csvautotabular{tables/essential_features.csv}

	\subsection{Desirable Features}
	\csvautotabular{tables/desirable_features.csv}

\end{landscape}

\section{Prototypes}
\subsection{GUI Prototype}
\subsection{Animation Prototype}

\chapter{Project Plan}

\bibliographystyle{IEEEtran}
\bibliography{interim}
\addcontentsline{toc}{chapter}{Bibliography}

\appendix
\chapter{Prototype Source Code}
\section{GUI Prototype}
\section{Animation Prototype}

\end{document}
